\section{Claim ledger and remaining work}
\label{sec:status}

\subsection{What is established in this report}

The following results are supplied with complete human proofs or proof
reductions:

\begin{enumerate}
\item the fibered finite-query interface and realizer definitions;
\item local certificate sealing through existential elimination;
\item composition through substitution and simultaneous cut;
\item the fixed-formula Boolean reflection construction, with explicit
      quantifier-free induction accounting;
\item bounded least accepted-certificate selection from the checked search
      theorem;
\item master-certificate formation for finite bounded decidable families;
\item least finite monotone demand closure with a sharp potential bound;
\item realization of fixed bounded strategy certificates;
\item the FIBER compilation theorem by induction over the source schedule;
\item the residual-elimination hierarchy;
\item the dyadic adjacent-motion/Cauchy theorem;
\item finite-net lifting with local witness sealing;
\item guarded bisection and exact branch uniqueness from a positive factor
      bound;
\item the finite-use impossibility of unguarded exact sign selection;
\item the diagonal impossibility of a primitive-recursive universal evaluator;
\item stable cluster, finite-algebra, and representative-relation interfaces.
\end{enumerate}

The exact reference algorithms included in the archive test mixed-radix
bundling, least selection, demand closure, strategy trees, dyadic bounds,
finite-net lifting, and certificate escape behavior.

\subsection{What is inherited, not reproved}

The report relies on checked repository declarations for the object calculus,
substitution, weakening, simultaneous cut, quantifier-free induction,
arithmetic, order, search, pairing, bounded sums, powers, integers, and several
linear-algebra constructions.  It also relies on the checked bridge for
transport to Foundation clients.  These dependencies are named in the source
audit.  They were inspected through the GitHub snapshot but were not rebuilt
inside this artifact environment.

\subsection{What is not yet kernel-checked}

No new Lean theorem in this archive was compiled.  In particular, the
following remain implementation milestones:

\begin{enumerate}
\item a complete object proof of the bounded conjunction/disjunction fold
      reflection theorems;
\item the external \texttt{BFm.reflect} proof generator;
\item the one-success adapter over \texttt{Core.Search};
\item a canonical sequence coding suitable for generated master certificates;
\item the generic master-certificate theorem;
\item fixed-arity and then variable-arity demand closure;
\item the bounded strategy compiler;
\item the typed graph compiler returning literal \texttt{Prf PRAAx} terms;
\item the integer absolute-value and scaling lemmas needed for the dyadic
      theorem;
\item finite-net and cluster examples in the object theory;
\item integration builds and axiom audits.
\end{enumerate}

The distributed Lean files therefore use comments and theorem targets for
unfinished arithmetic.  They contain no claim that a source file without
\texttt{sorry} has been elaborated.

\subsection{Research boundaries not removed by the framework}

FIBER resolves representational and compositional obstacles.  It does not by
itself prove:

\begin{itemize}
\item general effective Faltings, Siegel, Roth, or Mordell--Weil generation;
\item effective Chebotarev, Siegel--Walfisz, or Bombieri--Vinogradov constants
      beyond what is explicitly supplied;
\item a finite presentation or modulus for every abstract compact space;
\item a Cauchy rate from metastability alone;
\item an exact branch, sign, eigenvector, or axis selector at degeneracy;
\item rational representatives for arbitrary abstract descent classes;
\item a universal primitive-recursive evaluator;
\item or automatic extraction from every classical proof.
\end{itemize}

For each boundary the framework gives either a corrected theorem or an exact
producer formula.

\subsection{Recommended first integration slice}

A realistic first pull request should be intentionally small:

\begin{enumerate}
\item add \texttt{FIBER/Kernel.lean} and compile its structural aliases;
\item add the one-success search adapter and a two-query master certificate;
\item prove a fixed two-coordinate demand closure theorem;
\item prove the dyadic adjacent-motion theorem for one explicit signed-integer
      encoding;
\item compile a miniature example that consumes two independently certified
      approximants and exports a Cauchy estimate;
\item run \texttt{\#print axioms}, native transfer, and an open-sequent bridge
      check.
\end{enumerate}

This slice exercises the distinctive contribution of the framework without
first requiring a generic formula compiler.  Once it is checked, the bounded
formula and source-graph generators can be built against a stable kernel.

\subsection{Acceptance criteria for a packet migrated to FIBER}

A packet should count as migrated only when:

\begin{enumerate}
\item its main target is stated as an explicit ordinary \(\ipra\) formula;
\item every primitive symbol and representation convention is fixed;
\item the finite query and demand map are explicit;
\item every checker has an object soundness/reflection derivation;
\item every induction predicate is displayed and quantifier-free;
\item all certificate scopes are classified and enforced;
\item all missing data appear as residual formulas;
\item the endpoint is a compiled \texttt{Prf PRAAx} declaration;
\item its axiom audit and build command are recorded;
\item and any Foundation or application endpoint is obtained through a checked
      bridge or an independently stated theorem.
\end{enumerate}
